\scanpage{3}
\reversemarginpar
12)\margin{\unsure{commencé} par Raynaud}% slanted
\hspace{.4em}Mise au point, et extensions diverses, de la théorie des modèles de Néron. \hfill
\begin{minipage}[c]{1.8em}\figure{\tikz[baseline=-.35ex,x=1ex,y=1ex]{\draw (.1,2) -- (0,-.5); \draw (.36,1.85) -- (.24,-.45);}}\end{minipage}

\textcircled{\small 13}\margin{S'impose moins depuis les résultats d'approximation d'Artin}% slanted
\hspace{.4em}\unsure{Théorémes|Théorèmes} \unsure{assymptotiques|asymptotiques} pour les \unsure{foncteurs} de Greenberg, relatifs à un anneau local \unsure{noeth} complet ou hensélien.

14) Structures des Topos (questions d'existence de foncteurs fibres....). Structure des champs champs représentables.... \hfill
\begin{minipage}[c]{1.8em}\figure{\tikz[baseline=-.4ex,x=1ex,y=1ex]{\draw[double,double distance=.3pt] (.7,2) -- (1,2.1) -- (.46,-1) -- (.13,-.65);}}\end{minipage}

\textcircled{\small 15}\margin{fait partiellement par \unsure{Mazur}-Artin. Réflexions de Quillen \unsure{en} train}% slanted
\hspace{.4em}\unsure{Théorie|Théories} du \deleted{\unsure{g}}type d'homotopie des Topos, incluant une définition des $\pi_{\unsure{i|\bullet|\cdot}}$ à \deleted{\unsure{coefficients}} \unsure{noveaux|noyaux|nouveaux} infinitésimaux. \inserted{\unsure{Th's|Thés} d'\unsure{Hurwisz|Hurwicz}. Suites exactes d'homotopie. \unsure{Ths|Thés} de comparaison, \deleted{de \unsure{spécialisation}.}}

\textcircled{\small 16}\margin{contre-exemple de Artin, \unsure{situation} peu engageante}% slanted
\hspace{.4em}\deleted{Théorie de spécialisation du} \inserted{``vrai''} \deleted{groupe fondamental, pour une famille de courbes.} \hfill
\begin{minipage}[c]{1.8em}\figure{\tikz[baseline=-.4ex,x=1ex,y=1ex]{\draw (.75,1.65) -- (1,1.65) -- (.67,-.8) -- (.48,-.65);}}\end{minipage}

17)\margin{\unsure{résolu} par Raynaud \(\bigl[\)}% slanted
\hspace{.4em}\deleted{Séparation de foncteurs : classes de correspondance.... Cas de représentabilité et de non représentabilité de \mbox{$\underline{\mathrm{Pic}}_{X/S}$}, \mbox{$X$} de dim} \deleted{\unsure{de|rel}} \deleted{relative \mbox{$1$} sur \mbox{$S$}.}

\textcircled{\small 18} Question du \unsure{passage} au quotient dans cas artinien et de l'effectivité de données de descente plates. \deleted{\unsure{Construction|Constructions} \unsure{générale|générales} de corps de \unsure{modules|variables}.} \hfill
\begin{minipage}[c]{1.8em}\figure{\tikz[baseline=-.4ex,x=1ex,y=1ex]{\draw[double,double distance=.25pt] (.9,1.8) .. controls (1,1.4) and (.7,0) .. (.31,-.8) -- (.08,-.5);}}\end{minipage}

\tikz[baseline=(num.base)]{\node[draw,circle,inner sep=.7pt] (num) {\shortstack{19\\[-.5ex]\scriptsize$\unsure{\times|*}$}};}
Théorie d'\unsure{amplitude|amplitudes} pour cycles de dim. quelconque.
% archive sheet number: 2
\normalmarginpar

